\section{Discussion}
\label{sec:discussion}

This chapter measures the distance between the objective of \cref{sec:intro-objective} and the
results of \cref{sec:results}. It then answers the three hypotheses of \cref{sec:bg-gap}, and ranks
the changes that would close the gap. The ranking is the chapter's deliverable, and every entry
carries the measurement that ranks it.

\subsection{The distance to the objective}

The target was a bipedal 100\,m record pace near 4\,m/s. The measured mechanism sustains about
1.3\,m/s thermally, and bursts to about 2.7\,m/s (\cref{tab:speed-ceiling}). A 100\,m dash at
record pace is about 25\,s of continuous running against a continuous torque rating. The sustained
figure therefore decides the record, and the gap is a factor of three.

That gap is mechanical rather than algorithmic, and the evidence is a controlled comparison rather
than an inference. A 21-gain scripted controller with no learning in it hits the identical
pitch-balance wall at the identical time. The survival time equals the free-topple time of an
uncontrolled inverted pendulum, across a 5$\times$ change in the plant's time constant
(\cref{tab:freefall}). At that wall the set of stabilising action sequences is essentially empty,
so every controller drawn from it performs identically. No algorithm closes the gap on this plant.

\subsection{Verdict on the three hypotheses}

\textbf{H1, dimensionality reduction: supported.} The structural prior turned a non-converging
problem into a trainable one. The torque attempt is the control. It never converged, across an
unbroken sequence of reward-shaping attempts (\cref{sec:rl-torque}). The per-step position policy
converged only to standing, and skated (\cref{sec:rl-position}). The Fourier action space walked at
0.5\,m/s on its first trained policy, with no skating signature. The claim is scoped to this plant
and this sample budget. No arm was run at the $10^8$ to $10^9$ transitions per condition that the
published comparisons use.

\textbf{H2, control authority: partially supported, and the interesting half is untested.} Gait
morphology alone was not sufficient. Removing the per-step residual channel from the oscillator
generator produced stability without locomotion, at 15\,m of a 72\,m course (\cref{sec:cpg-ab}).
The reflex and residual channels were therefore necessary. The impedance level of the hypothesis,
phase-dependent $K_p(\varphi)$ and $K_d(\varphi)$, was designed and never run. The drive mode it
needs was never switched on the hardware (\cref{sec:rl-impedance}). H2 is answered for the residual
channel, and open for the impedance channel.

\textbf{H3, robustness across the sim-to-real boundary: untested on hardware.} No policy has run on
the robot. The closest available evidence is the descent onto the measured 0.8\,Hz drive. The
structured controller completed it, to an episode length of 4938 and three times its parent's score
(\cref{sec:sim2real-lineage}). That is a sim-to-sim result on a measured plant. No unconstrained arm
was run through the same descent for comparison, so it does not answer H3 as stated.

\subsection{Ranked levers for a successor robot}
\label{sec:future}

\begin{enumerate}[leftmargin=*]
  \item \textbf{Bus voltage before torque.} 87\,\% of feasible gaits are limited at the cam's
        peak-power corner (\cref{sec:speed-ceiling}). Over that population, speed scales with
        no-load speed and is insensitive to peak torque, and no-load speed is linear in bus
        voltage. The scope is the sweep as run, on the CAD plant. The cost is pack mass plus a
        drive rated for the higher voltage. A 60\,V variant of a drive already surveyed would have
        given 25\,\% for EUR\,220 (\cref{sec:hw-motor}).
  \item \textbf{Distal mass.} At $-0.93$\,(m/s)/kg it carries nine times the leverage of torso
        mass (\cref{tab:speed-ceiling}). A carbon shin, plus a proximally mounted spring on a
        tendon, recovers about 0.5\,m/s. The cost is a composite part and a tendon route to
        design. The carbon ankle already executed the cheap half of this, at $-209$\,g per shin
        (\cref{sec:ankle-build}).
  \item \textbf{Re-proportion or replace the 4-bar.} Only 0.20 to 0.35\,m of a 1.06\,m geometric
        span is usable. The standing pose sits at 96\,\% of reach, where the toe-force ellipse is
        19:1 anisotropic in the wrong direction (\cref{sec:hw-force-map}). Backward reach from the
        stance pose is 2.7\,cm against 31\,cm forward, which is why a backward fall can never be
        caught (\cref{sec:hw-experiments}). Target a flat transmission ratio over about 0.5\,m,
        with stance at about 85\,\% of reach. The cost is a linkage redesign. In the meantime,
        5\,cm of crouch buys most of the backward reach for 2\,N$\cdot$m out of 144.5.
  \item \textbf{Size the ankle for 3.5 body weights.} The passive ankle delivers a median of
        0.42\,BW over the running stance region, which is 8$\times$ short, and it fails in positive
        mechanical feedback (\cref{sec:hw-force-map}). The cost is either a stiffer passive element
        with correct preload, or a seventh and eighth motor. \Cref{sec:rl-pitch-wall} shows the
        passive route suffices in simulation, if the stiffness is achievable.
  \item \textbf{Cool the cam motor.} The demanded 167\,N$\cdot$m sits against a 55\,N$\cdot$m
        continuous rating, and the sustained figure is what decides a record. Cooling extends how
        long the peak stays available. The gain is unquantified, because winding temperature has
        never been logged at the measured duty cycle. \Cref{sec:speed-ceiling} records that as an
        open item. This is therefore a lever with a named test rather than a ranked result. The
        cost is added mass and an airflow path at the joint.
  \item \textbf{Treat torque-mode control as an architectural requirement from day one.} It was
        abandoned twice here, once on convergence and once on assumed deployability. The measured
        0.8\,Hz drive loop points back to it (\cref{sec:plant-id-drive}). Unlike the levers above,
        this one carries no number. In my view it is a design opinion rather than a ranked result.
\end{enumerate}

Two items sit outside the ranking, because they are measurement gaps rather than design changes.
Pack voltage sag is unobservable from the drives, and plausibly worth 30\,\% of joint speed at the
ten-second current tier (\cref{sec:hw-battery}). A voltmeter across the pack during a sprint
settles it. And the actuator's reference frame for reported current is undocumented, which puts an
unquantified risk on every measured torque figure (\cref{sec:hw-motor}). A clamp meter on a stalled
motor settles that one.

\subsection{Levers for this robot, in the order that unblocks the next}

\begin{enumerate}[leftmargin=*]
  \item \textbf{ONNX export and a Pi inference loop.} Everything else waits on it. The toy problem
        of \cref{sec:toy-problem} already demonstrated the path at one degree of freedom.
  \item \textbf{Current-mode control with the position PD closed on the Pi at 200\,Hz}, replacing
        the 0.8\,Hz drive loop that every measurement identifies as binding.
        \Cref{tab:loop-timing} shows the budget exists.
  \item \textbf{The carbon ankle strut}, already fitted, and the corresponding retrain on the
        rebuilt plant.
  \item \textbf{A from-scratch retrain with cadence pressure from step zero.} Every warm-started
        run stays trapped in the fast-stepping basin (\cref{sec:rl-cadence}).
  \item \textbf{The phase-dependent impedance experiment} of \cref{sec:rl-impedance}. Items 1 and 2
        make it possible, and it closes the open half of H2.
\end{enumerate}

\subsection{Method-level lessons}

The lessons below are the ones that survive a plant change, by the test of \cref{tab:survives}.

\textbf{Optimisers exploit objectives, and neural networks are not the culprit.} The policy gamed
the reward by skating, penetrating the floor, shuttling and committing suicide by penalty. The
scripted baseline's coordinate-descent tuner gamed its own objective identically. It drove at
1.6\,m/s against a 0.3\,m/s command, to farm a distance term (\cref{sec:scripted-baseline}). On
those two cases, in my view every objective needs adversarial review before its optimiser gets it.

\textbf{The margin advantage of learning rests on three premises, and none was paid in full here.}
The first is an exact plant, which this was not, since every parameter moved. The second is a
reward pricing distance from the failure boundary, which this did not, since it priced survival and
distance. The third is an asymptotic budget, which 600\,M steps per milestone is not. The
generalisable diagnostic is that seed variance is the margin measurement. A figure of
20.3\,$\pm$\,21.6\,s against 5 of 5 at the cap is bimodality rather than a worse mean, and a curve
plotting the mean will not show it.

\textbf{A gait prior reads on two axes, and the useful question is which kind rather than how
much.} A shape prior fixes the trajectory and exposes its timing, so tuning does not enlarge the
reachable set. It is therefore safe only when the morphology has a natural template. A structural
prior constrains the class and leaves the member free. DASH-01 separates the two, because its cam
turned out to be a full crank with no biological analogue. The evidence that the wider space was
genuinely used is the cadence agreement between the learned optimum and the independent physics
sweep, both at about 12.4\,steps/s (\cref{sec:speed-ceiling}). In my view the two-axis reading is a
way of comparing systems usually filed under one label, rather than a measurement.

\textbf{Every assumed parameter measured later moved pessimistically.} Mass by $+18\,\%$, torque by
$-15\,\%$, drive bandwidth by $-16\times$, plus friction and the ankle. Measure before training.
And when a measurement lands, apply it as \emph{the} plant rather than as an optional variant, or
the wrong default gets trained on (\cref{sec:plant-id}).

\textbf{Simulation results need the same scepticism as rewards.} Six artefacts each produced a
clean, converged, wrong result. They were soft joint limits, constraints inactive at settle,
frame-relative safety checks, a stale keyframe, a solver-version split between the development
machine and the cluster, and an actuator velocity cap that was never present. The last let a policy
balance at 23\,rad/s on a 22\,rad/s motor. Self-tests, smoke gates or measurement caught each one.

\textbf{Instrumentation found the ceilings, and cost something to build.} The measurement
instrument found the drive ceiling. The flight recorder found a daemon-killing crash and a dead bus
within its first hour (\cref{sec:deploy-blackbox}). The cost is that the measurement disturbed the
measured twice: once with the inertial sensor on a swaying rig, and once with the loop rate under
SSH load. The second produced a committed wrong conclusion that only an A/B retracted.

\textbf{Progressive validation localised the failures.} The base-DOF ladder localised failures to
single degrees of freedom. The toy problem debugged the deploy chain at one degree of freedom. Both
are cheap, and both worked.

\subsection{What a successor project should do differently}

The validated toe-force map and the capability sweep together make a \emph{design search in
simulation} the natural next use of the training cluster. Such a search would range over bus
voltage, gear ratio, distal mass, linkage proportions and ankle stiffness. It would search over
robots rather than over controllers, for a robot already known to be the constraint. The tooling
exists in validated form. The analysis runs on the same MJCF the policies train against, its
Jacobians are self-tested, and its predictions agree with MuJoCo within 1\,\%.

The sequencing lesson is narrower and cheaper. The measurement campaign that produced
\cref{tab:plant-id} took three days, and it reclassified the project. It ran in month five of nine.
Run in month one, it would have changed the actuator purchase, the ankle design and the choice of
action space. Roughly 3 billion environment steps would then have been spent on the robot that
exists.
